\hypertarget{numerical-high-performance-data-structures}{%
\chapter{NUMERICAL HIGH-PERFORMANCE DATA
STRUCTURES}\label{numerical-high-performance-data-structures}}

\hypertarget{numpy-contiguous-layouts-and-strides}{%
\subsection{28.1 NumPy Contiguous Layouts and
Strides}\label{numpy-contiguous-layouts-and-strides}}

NumPy arrays are designed to speed up mathematical operations on
homogeneous datasets by utilizing contiguous blocks of memory. Unlike
standard Python lists, which contain pointers to scattered object
instances, a NumPy array stores the raw values directly in a single,
contiguous array block, which maximizes cache hit rates.

\hypertarget{c-contiguous-vs.-fortran-contiguous-layouts}{%
\subsubsection{1. C-Contiguous vs.~Fortran-Contiguous
Layouts}\label{c-contiguous-vs.-fortran-contiguous-layouts}}

The mapping of multi-dimensional arrays to linear virtual memory is
defined by the array's layout: * \textbf{C-Contiguous (Row-Major)}:
Elements along rows are stored adjacent to each other in memory. The
last index changes fastest when traversing the memory sequentially. This
is the default layout in NumPy. * \textbf{Fortran-Contiguous
(Column-Major)}: Elements along columns are stored adjacent to each
other in memory. The first index changes fastest when traversing memory
sequentially.

\begin{lstlisting}
2D Array Layout (3x3):
  [[1, 2, 3],
   [4, 5, 6],
   [7, 8, 9]]

C-Contiguous (Row-Major):
  [ 1 | 2 | 3 | 4 | 5 | 6 | 7 | 8 | 9 ]

Fortran-Contiguous (Column-Major):
  [ 1 | 4 | 7 | 2 | 5 | 8 | 3 | 6 | 9 ]
\end{lstlisting}

\hypertarget{stride-calculations}{%
\subsubsection{2. Stride Calculations}\label{stride-calculations}}

NumPy locates any element in an N-dimensional array using its
\textbf{strides} vector. A stride defines the number of bytes that must
be skipped in memory to move to the next index in a given dimension.

For a 2D array of shape \((R, C)\) and data type size \(D\) (in bytes),
the stride vectors are:

\[\text{Strides}_C = (C \times D, D) \quad (\text{Row-Major / C-Contiguous})\]

\[\text{Strides}_F = (D, R \times D) \quad (\text{Column-Major / Fortran-Contiguous})\]

\[\text{Memory Address} = \text{Base Address} + \sum_{i=0}^{N-1} (\text{Index}_i \times \text{Stride}_i)\]

For example, given a 3x3 array of 64-bit floats (\(D = 8\) bytes) in
C-contiguous layout: * \passthrough{\lstinline!strides!} =
\passthrough{\lstinline!(24, 8)!}. To move to the next row, jump 24
bytes (3 elements). To move to the next column, jump 8 bytes (1
element).

Accessing elements along strides that do not match the contiguous memory
direction (e.g.~slicing rows from a column-major array) breaks CPU cache
locality, as the processor must retrieve values from scattered memory
addresses (cache misses).

\hypertarget{vectorization-and-simd-pipelines}{%
\subsubsection{3. Vectorization and SIMD
Pipelines}\label{vectorization-and-simd-pipelines}}

Contiguous arrays enable compiler vectorization. Instead of executing
operations on individual elements sequentially (Scalar execution),
modern CPU architectures support \textbf{Single Instruction, Multiple
Data (SIMD)} instructions (e.g.~AVX-512, ARM NEON).

\begin{lstlisting}
Scalar Addition (Iterative loops):
  Step 1: A[0] + B[0] -> C[0]
  Step 2: A[1] + B[1] -> C[1]
  ...

SIMD Vector Addition (Single clock cycle):
  Vector Register A: [ A[0] | A[1] | A[2] | A[3] ]
                               +
  Vector Register B: [ B[0] | B[1] | B[2] | B[3] ]
                               =
  Vector Register C: [ C[0] | C[1] | C[2] | C[3] ]
\end{lstlisting}

By storing numbers contiguously, CPython's underlying math libraries
(like BLAS or LAPACK) load entire memory blocks into hardware vector
registers, computing operations across multiple elements in a single CPU
clock cycle.

\begin{center}\rule{0.5\linewidth}{0.5pt}\end{center}

\hypertarget{pyarrow-columnar-format-and-zero-copy-sharing}{%
\subsection{28.2 PyArrow Columnar Format and Zero-Copy
Sharing}\label{pyarrow-columnar-format-and-zero-copy-sharing}}

For large-scale data analysis or microservice architectures, exchanging
dataframes between distinct processes (like C++ loaders, Python
analytics scripts, and JVM processing nodes) introduces performance
bottlenecks.

\hypertarget{apache-arrow-columnar-format}{%
\subsubsection{1. Apache Arrow Columnar
Format}\label{apache-arrow-columnar-format}}

\textbf{PyArrow} implements the Apache Arrow format, which defines a
standardized, language-independent \textbf{columnar memory layout}.

\begin{lstlisting}
Record-Oriented Layout (Standard CSV / Database):
[Row 0: ID, Name, Salary] | [Row 1: ID, Name, Salary] | ...

Columnar-Oriented Layout (Apache Arrow):
[All IDs (contiguous)] | [All Names (contiguous)] | [All Salaries (contiguous)]
\end{lstlisting}

Columnar storage optimizes performance in three ways: *
\textbf{Vectorized Scan Operations}: The CPU can scan a single attribute
column (e.g.~calculating average salary) without loading unrelated
columns (like Name or ID) into cache memory. * \textbf{Contiguity}:
Column values are stored contiguously, allowing direct SIMD instruction
execution. * \textbf{Null Bitmaps}: Null values are tracked in a
separate bit mask, avoiding sentinel values inside the data array.

\hypertarget{zero-copy-sharing-protocols}{%
\subsubsection{2. Zero-Copy Sharing
Protocols}\label{zero-copy-sharing-protocols}}

Because Arrow memory layouts are identical across languages, PyArrow can
map memory buffers across process boundaries using shared memory or
plasma stores without performing copying or serialization.

Below is an example mapping data between PyArrow tables and Python
memory layouts with zero-copy overhead:

\begin{lstlisting}[language=Python]
import pyarrow as pa
import numpy as np

# Create a contiguous NumPy array
np_data = np.arange(1_000_000, dtype=np.int64)

# Wrap the NumPy array directly in a PyArrow Buffer
# PyArrow shares the memory address pointer of the NumPy array (zero-copy)
arrow_array = pa.Array.from_numpy_to_arrow(np_data)

# Create an Arrow Table
table = pa.Table.from_arrays([arrow_array], names=['metrics'])

# Verify that the underlying memory address remains identical
np_address = np_data.__array_interface__['data'][0]
arrow_address = table.column(0).chunks[0].buffers()[1].address

print(f"NumPy Memory Address: {hex(np_address)}")
print(f"Arrow Memory Address: {hex(arrow_address)}")
print(f"Are addresses identical (Zero-Copy)? {np_address == arrow_address}")
\end{lstlisting}

--- Papermill, pandas, and machine-learning frameworks (like PyTorch or
TensorFlow) utilize PyArrow columns as direct input sources, bypassing
serialization bottlenecks and memory duplication.

\begin{center}\rule{0.5\linewidth}{0.5pt}\end{center}
